\section{Ampliación del corpus y validación externa}
\label{sec:validacion-externa}

Las cifras de la Sección~\ref{sec:pilarA-evaluacion} tienen dos límites: los patrones se diseñaron
mirando el mismo corpus con que se evalúan, y ese corpus no contiene actividad legítima de fondo.
Esta sección los aborda con conjuntos que no intervinieron en el diseño. Cada uno se carga con sus
propios \textit{topics} y grupos de consumo, y se comprueba que los resultados de los 110 conjuntos
del Pilar~A no cambian (Anexo~\ref{anx:validacion}). \textbf{Las cifras oficiales de la fase no se
tocan}.

\subsection{Segunda versión de los patrones}
\label{sec:pilarA-v2}

La segunda versión (v2) corrige tres causas de error encontradas en la revisión:
\begin{itemize}
    \item P1 descarta las alertas de credenciales cuyo origen es el ruido de arranque
    (\texttt{csrss.exe} y \texttt{wininit.exe});
    \item P2 solo acepta el controlador de dominio como destino si hay ejecución remota desde una
    herramienta de movimiento lateral;
    \item el evaluador Sigma interpreta las palabras clave como subcadenas, como prevé Sigma, lo
    que recupera 1.692 alertas de dos reglas.
\end{itemize}
Sobre las capturas atómicas, la precisión de P1 pasa de 0,43 a 0,56 y su exhaustividad de 0,47 a
0,74, en evaluación automática. Recupera Mimikatz ejecutado en memoria y quita tres campañas
falsas. Como el ajuste se hizo mirando esos datos, se portó a Flink como versión paralela, con
salida propia, y \textbf{solo se adopta tras validarla en conjuntos que ningún ajuste ha visto}.

\subsection{Conjuntos y resultados}

\begin{table}[H]
\centering
\small
\setlength{\tabcolsep}{4pt}
\begin{tabular}{|p{2.6cm}|p{3.4cm}|r|r|r|p{3.3cm}|}
\hline
\textbf{Conjunto} & \textbf{Qué aporta} & \textbf{Eventos} & \textbf{Alertas} & \textbf{Camp.} & \textbf{Lectura} \\
\hline
APT29, segundo día \cite{otrf} & Otra jornada de ataque sobre la misma red & 585.479 & 290 & 5 / 6 & v1: 4 de 5 verdaderas; v2: todas verdaderas \\
Zeek de APT29 \cite{otrf} & Primera fuente de red & 2.088 & 0 & 0 & Es de otra ejecución: solo topología \\
COMISET, red real \cite{comiset} & Actividad normal de 15 equipos (50,6\,h) & 1,67\,M & 48 & 0 & 0,95 alertas y 0 campañas por equipo y hora \\
COMISET, equipo con pruebas & Atomic Red Team \cite{atomicredteam} sin etiquetar & --- & 1.183 & 4 & Ataque simulado reconocido \\
COMISET, laboratorio \cite{comiset} & Un equipo con ataques etiquetados por reglas & 600.218 & 85 & 0 & v1 = v2; etiquetas ruidosas \\
Laboratorio propio & Verdad por ventanas temporales & 14.945 & 886 & 1 & P1 correcto; 0 en la línea base \\
\hline
\end{tabular}
\caption{Conjuntos de la validación externa. En el segundo día de APT29 se indican las campañas de
v1 y de v2; en los demás, las dos versiones coinciden.}
\label{tab:validacion-externa}
\end{table}

\medskip\noindent\textbf{APT29, segundo día y Zeek.} En el segundo día, v1 da dos campañas verdaderas de P1
con Mimikatz, una de P2 y una de P3, y una falsa de P1 con el ruido de \texttt{wininit.exe}. En v2,
la falsa desaparece, aparece la falsificación de un \textit{Golden Ticket} y P2 recupera el salto al
controlador de dominio: \textbf{todas sus campañas son verdaderas}. Zeek aporta la primera fuente de
red, con su mapeador a OCSF, pero no confirma ninguna arista de P2: sus registros son 25 y 51 horas
anteriores a los de los equipos, de otra ejecución de la emulación. Sirven para la topología, no para
correlacionar.

\medskip\noindent\textbf{Falsas alarmas sobre actividad normal.} COMISET \cite{comiset} publica 14 días de
una red universitaria real: 16 equipos y 202 millones de registros. Tras un censo completo se
eligieron 8 horas por equipo, antes de detectar nada. Sobre 50,6 horas observadas de 15 equipos
salen \textbf{48 alertas, 3 de nivel alto y ninguna campaña}. Son 0,95 alertas por equipo y hora, con
una cota superior del 95\,\% de 0,06 campañas por equipo y hora. Las alertas son mantenimiento de
Windows. Es la tasa que las capturas atómicas no podían dar. El censo encontró además un equipo con
Atomic Red Team ejecutando pruebas: se separó de la medida, y en él el sistema abre \textbf{cuatro
campañas} de P1 y P3 sobre una actividad que el conjunto no etiqueta.

\medskip\noindent\textbf{Laboratorio de COMISET.} Es un solo equipo con eventos etiquetados con ATT\&CK. Se
tomaron las 12 horas con más etiquetas de las técnicas de los patrones y 6 sin ninguna como
control. \textbf{v1 y v2 dan las mismas 85 alertas, ninguna campaña y ninguna alerta en el control.}
Las etiquetas no son una verdad independiente: las pone un sistema de reglas, y la mayoría son
comportamiento normal de Windows. Sí hay un ataque real: un servicio de control remoto vulnerable
lanza una consola como SYSTEM y le siguen reconocimiento y un PowerShell interactivo como SYSTEM.
El sistema lo alerta en el momento exacto, pero \textbf{P3 no abre campaña}, porque ninguna regla
cubre la explotación del servicio y el patrón necesita esa etapa.

\medskip\noindent\textbf{Laboratorio propio.} El conjunto de la Sección~\ref{sec:laboratorio} es \textbf{la
única verdad independiente de cualquier regla}, porque sale de las ventanas del manifiesto. Se carga
la telemetría completa, separada en el ataque (1\,h 12\,min) y la línea base (5\,h 49\,min). P1
reconoce el robo de credenciales con 19 evidencias: \texttt{procdump} y \texttt{comsvcs.dll} sobre
LSASS, la limpieza con \texttt{fsutil} y la herramienta en disco. La línea base no produce ninguna
campaña. De las cuatro técnicas, tres tienen alertas específicas en su ventana; la inyección, que
Defender bloqueó, no tiene ninguna.

\subsection{Decisión y límites}

En los dos conjuntos que ningún ajuste ha visto, \textbf{v2 produce las mismas alertas y campañas que
v1}, alerta por alerta. En las capturas atómicas y en el segundo día de APT29, en cambio, corrige
errores identificados. Por eso \textbf{v2 pasa a ser la versión actual} para las fases siguientes y
el portal. Las cifras v1 se mantienen como las de referencia de la fase. Hay tres límites:
\begin{itemize}
    \item la actividad normal de COMISET no está etiquetada y se supone benigna;
    \item los conjuntos externos no permiten medir la exhaustividad;
    \item los dos falsos negativos encontrados son de cobertura de reglas, no de la lógica de los
    patrones.
\end{itemize}
